\section{Completing the computation}\label{sec:assembly}

We now have the two certificate families described in
Section~\ref{sec:certificates}. Let $\alpha$ be the threshold constructed
in Theorem~\ref{cp:threshold}. Proposition~\ref{cp:finite-computation}
gives a computable rational sequence $(\ell_n)$ with
$\ell_n<\alpha$ and $\ell_n\to\alpha$. Its underlying finite-size bound
is the exact one-sided estimate
\[
 \frac{f_n}{n}-2^{67}n^{-1/6}\le\alpha.
\]
For the upper certificates, a rational trial is the finite site-only
description of Section~\ref{ctr:section}: it has strictly increasing
rational exponents in $(0,1)$ and a rational-valued message pair on a
finite rational rectangular grid of the root and site uniform variables.

\begin{lemma}[Upper certificates]\label{assembly:upper}
For positive rational $a$, positive integer $\beta$, and a rational trial
$Q$, the inequality $G_1(a,\beta;Q)<0$ certifies $\alpha\le a$.
Every rational $a>\alpha$ has such a certificate. The class of certified
densities is effectively enumerable.
\end{lemma}

\begin{proof}
The trial bound of Proposition~\ref{ctr:upper} gives
$P(a,\beta)\le G_1(a,\beta;Q)$. If $a<\alpha$, the pressure exists by
Lemma~\ref{ctr:pressure} and is nonnegative by Corollary~\ref{cp:soft-signs}.
Thus a negative trial is impossible below $\alpha$ and certifies the
non-strict inequality $\alpha\le a$. Nothing about the pressure at
equality is needed.

If $a>\alpha$, Corollary~\ref{cp:soft-signs} supplies a positive integer
$\beta$ for which $P(a,\beta)<0$. Theorem~\ref{cc:complete} says that
$P(a,\beta)$ is the infimum of the rational trial values. A rational
trial therefore has negative value. Its finite description is eventually
enumerated, and Lemma~\ref{ctr:evaluation} eventually gives a rational
interval for its value with negative upper endpoint. Enumerating positive
rational densities, positive integer penalties, finite descriptions and
accuracy indices gives the asserted enumeration.
\end{proof}

\begin{proof}[Proof of Theorem~\ref{thm:main}]
Theorem~\ref{cp:threshold} proves the strict-side probability limits for
the proper three-clause model in the introduction, with
$\alpha\in[1/200,10]$. Set $\alpha_3=\alpha$. This also identifies the
constant with the same-model $k=3$ threshold in
\cite[Theorem~1.1]{FixedClauseThreshold} by the uniqueness argument in
Section~\ref{sec:proof-map}.

Use the sequence from Proposition~\ref{cp:finite-computation} as the lower
sequence in Lemma~\ref{cert:search}, and use $G=G_1$ on the rational trial
class. This class is effectively enumerable by its finite encoding, and
Lemma~\ref{ctr:evaluation} proves uniform computability of its values.
Lemma~\ref{assembly:upper} proves both sign hypotheses of the search
lemma. Since $\alpha\in[0,20]$, all its assumptions hold. The resulting
single finite machine halts on every unary input $1^r$ and returns the
required binary-encoded rational within $2^{-r}$.

The compactness subsequences, conditional laws, Gaussian parameters and
approximation tolerances used to prove completeness are not inputs to the
machine. They establish the existence of finite witnesses, which the
search finds through its fixed enumeration. The running time of this fixed
machine on $1^r$ is itself a total computable function of $r$, obtained by
simulating the halting computation; the proof supplies no efficiency
estimate for that function.
\end{proof}
